\section{Conclusion}
\label{sec:conclusion}

This paper investigated to what extent LLMs can replace human involvement in configuring end-to-end data integration pipelines. Across three case studies, we compared an LLM-based pipeline against human-configured baselines.

Overall, the LLM pipeline produced similar results as the human-configured pipelines for all three case studies. This was confirmed by the step-wise as well as the structural end-to-end evaluation.  

LLM-based schema matching achieved perfect accuracy (67 correct out of 67 correspondences), including datasets with non-descriptive column headers where the LLM inferred semantics from values alone. For entity matching, LLM-labeled training data produced matchers of comparable quality to human-labeled data (0.937 vs.\ 0.916 average F1); since both pipelines share the same training code, this comparison isolates label quality directly. For data fusion, RAG validation sets achieve 77.3\% average test accuracy compared to 81.6\% for human-created ones, with the gap concentrated in the Companies use case (.744 vs.\ .861) where time-sensitive attributes cause systematic mismatches between outdated values in the datasets and the LLM's current knowledge. For domains with predominantly static attributes (Music, Games), the gap is negligible. Configuration time drops on average from an estimated 19 person-hours to 2 hours of unattended compute at approximately \$9 in LLM usage fees, while execution times remain similar (88\,s vs.\ 77\,s). The labor savings are substantial: For our use cases, the LLM pipeline eliminated the need for manual schema mapping, training data labeling, and fusion validation set creation entirely.

\textbf{Limitations.} Our pipeline assumes flat source tables with 1:1 attribute correspondences and does not cover scenarios where entity data is distributed across normalized tables, nor more complex correspondence types (1:n or higher-order correspondences, aggregation). All three case studies use publicly available datasets that could be contained in GPT-5.2's training data, so the strong results, particularly for schema matching and entity labeling, may partly reflect memorized knowledge rather than general reasoning. Performance on proprietary or domain-specific data outside the LLM's training distribution remains an open question. The temporal knowledge gap observed in Companies fusion points to a related limitation: LLM-generated validation data is most reliable for well-known entities with time-independent attributes, and less so for time-sensitive or obscure records. 

\textbf{Future Work.} 
Promising directions for future work include the agent-based refinement~\cite{zhu2025surveydataagents} of the automatic pipeline, as well as exploring the use of the pipeline’s outputs as a starting point for collaborative human–agent refinement workflows~\cite{santos2025interactiveHarmonization}. For value normalization, LLM-based coding agents could dynamically generate transformation functions that combine the robustness of code-based normalization with the LLM's ability to adapt to heterogeneous value formats. The pipeline currently runs strictly forward: each step consumes the output of the previous one without revisiting earlier decisions. An agentic architecture that feeds back downstream results to earlier steps could improve quality, for instance by detecting oversized entity clusters after matching and selectively relabeling the training pairs that contributed to false correspondences, or by identifying values of diverting format in the value clusters before fusion and using this information to refine the normalization step. Human-agent collaboration offers a middle ground where practitioners review and correct LLM decisions at critical points~\cite{santos2025interactiveHarmonization, Barbosa2026}. Today, the research on end-to-end data integration pipelines is hampered by the lack of challenging benchmarks that cover all steps of the integration process with ground truth. We publish the artifacts of our three case studies as an initial step toward such benchmarks.